\cleardoublepage
\section{Simplifying Expressions}

In this section, we'll learn how to combine like terms and use the distributive property to make expressions cleaner and easier to work with. Simplifying expressions helps make your work cleaner and your thinking clearer. It's like cleaning up a messy desk so you can find what you need.

\textbf{Like terms} have the same variable and exponent. You can add or subtract them just like regular numbers. The \textbf{distributive property} lets you spread a number across parentheses:
\[
a(b + c) = ab + ac
\]

\subsection*{Example: Combining Like Terms and Distribution}

Let's look at some examples of simplifying expressions:

\textbf{Combining Like Terms:}
\[
2x + 3x = 5x
\]

\textbf{Using the Distributive Property:}
\[
3(x + 4) = 3x + 12
\]

Whether you're solving equations or designing a budget, simplification makes everything easier. Even NASA engineers and video game developers simplify expressions when building models and writing code. It's a skill that goes far beyond school!

When simplifying, circle or highlight like terms before you combine them. Be careful with minus signs—they belong to the number or variable that follows! Go one step at a time to avoid careless mistakes.


\cleardoublepage
\subsection*{Short Question (English)}
Simplify: $7x + 5x$

\fullpageimage{images/q2_2_1-eng.jpg}

\newpage
\subsection*{Short Question (Korean)}
Simplify: $7x + 5x$

\fullpageimage{images/q2_2_1-kor.jpg}

\newpage
\subsection*{Short Question (Answer)}
Simplify: $7x + 5x$

\textbf{Answer:} $\mathbf{12x}$

\cleardoublepage
\subsection{Long Question (English)}
Use the distributive property to simplify $4(a + 6)$.

\fullpageimage{images/q2_2_2-eng.jpg}

\newpage
\subsection{Long Question (Korean)}
Use the distributive property to simplify $4(a + 6)$.

\fullpageimage{images/q2_2_2-kor.jpg}

\newpage
\subsection{Long Question (Answer)}
Use the distributive property to simplify $4(a + 6)$.

\textbf{Answer:} $4a + 24$

\cleardoublepage
\subsection*{Challenge Question (English)}
Simplify: $3x + 7 - 2x + 4$

\fullpageimage{images/q2_2_3-eng.jpg}

\newpage
\subsection*{Challenge Question (Korean)}
Simplify: $3x + 7 - 2x + 4$

\fullpageimage{images/q2_2_3-kor.jpg}

\newpage
\subsection*{Challenge Question (Answer)}
Simplify: $3x + 7 - 2x + 4$

\textbf{Answer:} Combine like terms: $(3x - 2x) + (7 + 4) = x + 11$

